\section{The Beginning of the World}

The Dominicans and Franciscans differed in their view of the age of the world. The Dominican \textbf{Thomas Aquinas} claimed that reason alone could not determine whether the world had a beginning in time or not. In this, he was followed by \textbf{Kant} who also famously claimed that pure reason could not resolve this question.

However, the Franciscan \textbf{Bonaventure} did demonstrate that the world has a beginning in time, while demonstrating a knowledge of the infinite centuries before its mathematical explanation was developed by \textbf{Cantor}. Bonaventure has three objections to a world with no beginning in time:

\begin{enumerate}


\item It is impossible to add to the infinite.

Mathematically, this is expressed by $\aleph + 1 = \aleph$.

That is, adding one minute to the infinite does not increase it. So, if we are at an infinite point of time from the beginning, it is not possible to add another minute.

\item It is impossible for the infinite to be ordered.

This is not completely correct, as the natural numbers are well ordered, symbolized by $\omega$. So you can go backwards by a finite amount without losing the order, e.g., $1 + \omega = \omega$, and so on. However, the negative integers, in their natural sequence, is not well ordered. Hence, it is impossible to count the days from the beginning.

\item It is impossible to traverse what is infinite.

That is, it would take an infinite number of days to reach our current era, which is impossible.
\end{enumerate}

Aquinas' objections are inadequate. Guenon has mentioned that the Middle Ages had an inadequate notion of the infinite, hence their metaphysics was incomplete. The example of Bonaventure shows this was not entirely true.

Note that Bonaventure's proof shows that \textbf{Nietzsche}'s doctrine of \textbf{Eternal Return} is false. This doctrine states that time is infinite, while space is not. Hence, the world must repeat itself identically over and over again. However, we now see that this is impossible, since we can never get to the world we are in.

Metaphysically, time, like space, is the realization of an idea or possibility. A world consists of a certain set of possibilities. When all the possibilities of a world have been manifested, that world ends. However, possibilities are infinite, hence there must be multiple worlds. These cannot be related to each other in time, due to Bonaventure's objections, but must be simultaneous in some sense, even if they appear to follow each other in time.

This again refutes Nietzsche's doctrine of a world with a finite number of possibilities which must repeat themselves. For Guenon, possibilities cannot repeat, since they are infinite. If they repeated, that would limit possibilities, which is a logical contradiction.

For background reading, please see \textit{Did the World Have a Beginning}\footnote{\url{http://www.philosophynow.org/issues/44/Did_the_World_Have_a_Beginning}}.


\flrightit{Posted on 2012-04-03 by Cologero}

